\input{preamble.tex}

\DeclareMathOperator{\im}{\operatorname{im}}

\title{\huge{MAT242 - Topology}}
\author{\LARGE{Thobias Høivik}}
\date{Fall 2026} 

\begin{document}
\maketitle

\newpage
\tableofcontents

\newpage
\section*{Intro}
The following are my notes for Topology which I will 
(hopefully) be taking in the fall of 2026. 
As these are my own personal notes I may skip ceirtain 
topics I already understand, be sloppy at times and 
this may cause trouble for anyone who wishes 
to use these notes for themselves so keep that in 
mind if you want to use these. 

\section{Topological spaces and continuity}
We recall from Real Analysis that for a metric 
space $(X,d)$ we have a notion of distance 
and can start describing open sets 
$U \subseteq X$ as sets where for every $x \in U$ 
there exists $\varepsilon > 0$ such that 
\[
    B(x;\varepsilon) \subseteq U, 
\]
meaning that for any $x$ in $U$ there is 
some open ball around it contained entirely in $U$.

In Hilbert spaces, normed spaces, $\mathbb R^n$ (different 
types and instances of metric spaces), etc., this is the usual idea: 
an open set is one where every point has some wiggle room around it.

Topology starts from the observation that many arguments in 
Analysis don't really use the distance $d(x,y)$. They only 
use collections of open sets. 

So instead of letting open sets be generated by some metric, why 
don't just specify which sets are open? 

Intuitively, a topology on a set $X$ is a choice about which 
sets are "open", whatever that may mean. 
Informally a topology tells us which points are "near" eachother 
without really specifing quantitatively what that may mean. 

So we want the least structure possible to be able to talk 
about things like continuity, connectedness, compactness, 
convergence, etc. 

\begin{defn}[Topological space]
    Let $X$ be a set. A \textbf{topology} 
    on $X$ is a collection $\tau \subseteq 
    \mathcal P(X)$ satisfying: 
    \begin{enumerate}
        \item $\emptyset, X \in \tau$.
        \item If $\{U_i\}_{i\in I}$ is any collection of elements 
            in $\tau$, then 
            \[
                \bigcup_{i\in I}U_i \in \tau. 
            \]
        \item If $U_1, \ldots, U_n \in \tau$, then 
            \[
                U_1 \cap \cdots \cap U_n \in \tau. 
            \]
    \end{enumerate}
    The pair $(X,\tau)$ is called a \textbf{topological space}. 
    The elements of $\tau$ are called \textbf{open sets}.
\end{defn}
We will often refer to a topological space by just 
the set if it is clear what the topology on $X$ is.  

The axioms are not arbitrary, but abstracted from metric spaces. 
In metric spaces we already have these properties even if they 
are not explicitly required. 
For instance in $\mathbb R$, 
\[
    \bigcap_{n=1}^\infty \left(-\frac{1}{n},\frac{1}{n}\right) 
    = \{0\}. 
\]
Each interval is open, but $\{0\}$ is not. 

Intuitively, finite intersections preserve some positive amount of 
wiggle room while infinite intersections can squeeze it all away. 

\subsection*{Examples of topological spaces}

\begin{ex}[$\mathbb R$]
    Let $X = \mathbb R$. The usual topology is the 
    collection of subsets $U \subseteq \mathbb R$ such 
    that for every $x \in U$, there exists $\varepsilon > 0$ with 
    \[
        (x - \varepsilon, x + \varepsilon) \subseteq U.
    \]
    This is the topology we know from Real analysis. 
    Examples of open sets include: 
    \[
        (0,1), \ (-3,7), \ \mathbb R, \ \emptyset, 
    \]
    and also things like 
    \[
        (0,1) \cup (2,5)\cup (10, \infty). 
    \]
    Non-examples: 
    \[
        [0,1], \ \mathbb Q, \ \{0\}. 
    \]
    Note that $\mathbb Q$ can be made into a topological space 
    by giving it the subspace topology from $\mathbb R$. 
    More on that later. 
\end{ex}
\begin{ex}[Metric spaces]
    Every metric space $(X,d)$ gives a topology. 
    Define $\tau_d$ by saying taht $U \subseteq X$ is open 
    if for every $x \in U$ there exists $\varepsilon > 0$ such that 
    \[
        B(x;\varepsilon) \subseteq U.
    \]
    Then $\tau_d$ is a topology on $X$. 
    So every metric spaces is a topological space, but 
    note that the opposite implication does not hold which 
    is one of the reasons topology is more general. 
\end{ex}
\begin{ex}[The discrete topology]
    Let $X$ be any set. Let $\tau := \mathcal P(X)$. 
    Then $\tau$ is a topology on $X$ which we call 
    the \textbf{discrete topology}. 
\end{ex}
\begin{ex}[The indiscrete topology]
    Let $X$ be any set. Define $\tau = \{\emptyset, X\}$. 
    This is called the \textbf{indiscrete topology} or 
    \textbf{trivial topology}.
\end{ex}

\begin{defn}[Closed set]
    A subset $C \subseteq X$ is called \textbf{closed} if its 
    complement is open: 
    \[
        C \ \text{is closed} \Leftrightarrow 
        X \setminus C \ \text{is open}.
    \]
\end{defn}

In a topological space, $\emptyset$ and $X$ are closed. 
Arbitrary intersections of closed sets are closed and 
finite unions of closed sets are closed. 

An important nuance is that closed does not mean not open. 
$\emptyset, X$ are both open and closed. In $\mathbb R$ (with the 
usual topology) we have that
$\emptyset, \mathbb R$ are both open and closed while e.g. 
$[0,1)$ is neither open nor closde. 

\subsection*{Continuity}
In Analysis, a function on metric spaces $f : X \to Y$ 
is continuous at $x\in X$ if small changes in $x$ in 
$f(x) \in Y$. Recall the $\varepsilon$-$\delta$ definition: 
\[
    \forall \varepsilon > 0, \ \exists \delta > 0 \ \text{s.t.} \ 
    d_X(x,y) < \delta \Rightarrow d_Y(f(x), f(y)) < \varepsilon.
\]

Since we don't have access to metrics we must rewrite this. 

\begin{defn}[Continuity]
    Let $(X,\tau_X)$ and $(Y,\tau_Y)$ be topological spaces. 
    A function 
    \[
        f : X \to Y
    \]
    is \textbf{continuous} if for every open set $V \in \tau_Y$, 
    the preimage 
    \[
        f^{-1}(V)
    \]
    is open in $X$. 
\end{defn}

\begin{prop}
    Let $(X,d_X)$ and $(Y,d_Y)$ be metric spaces. 
    A function 
    \[
        f : X \to Y
    \]
    is continuous in the $\varepsilon$-$\delta$ sense if 
    and only if for every open set $V \subseteq Y$, the 
    preimage 
    \[
        f^{-1}(V)
    \]
    is open in $X$.
\end{prop}
\begin{proof}
    First suppose $f$ is continuous in the $\varepsilon$-$\delta$ 
    sense. 
    Let $V \subseteq Y$ be open. We want to show that $f^{-1}(V)$ is 
    open in $X$. Take $x \in f^{-1}(V)$. Then 
    $f(x) \in V$. Since $V$ is open, there exists 
    $\varepsilon > 0$ such that 
    \[
        B(f(x) ; \varepsilon) \subseteq V. 
    \]
    By continuity of $f$ at $x$, there exists $\delta > 0$ such that 
    \[
        d_X(x,y) < \delta \Rightarrow d_Y(f(x),f(y)) < \varepsilon. 
    \]
    Thus, if $y \in B(x;\delta)$, then 
    \[
        f(y) \in B(f(x);\varepsilon) \subseteq V. 
    \]
    Therefore 
    \[
        y \in f^{-1}(V). 
    \]
    So 
    \[
        B(x;\delta) \subseteq f^{-1}(V). 
    \]
    Since every point of $f^{-1}(V)$ has a ball around it 
    contained in $f^{-1}(V)$, the set $f^{-1}(V)$ is open. 
    
    Now suppose that for every open set $V \subseteq Y$, the 
    preimage $f^{-1}(V)$ is open in $X$. 
    Let $x \in X$, and let $\varepsilon > 0$. The ball 
    \[
        B(f(x) ; \delta)
    \]
    is open in $Y$. Therefore 
    \[
        f^{-1}(B(f(x);\delta))
    \]
    is open in $Y$. 
    Since $x \in f^{-1}(B(f(x);\delta))$, there exists 
    $\delta > 0$ such that 
    \[
        B(x;\delta) \subseteq f^{-1}(B(f(x);\delta)). 
    \]
    Thus, if 
    \[
        d_X(x,y) < \delta, 
    \]
    then 
    \[
        y \in f^{-1}(B(f(x);\delta)), 
    \]
    so 
    \[
        f(y) \in B(f(x);\delta). 
    \]
    Equivalently, 
    \[
        d_Y(f(x),f(y)) < \varepsilon. 
    \]
    Hence $f$ is continuous at $x$. Since $x$ was arbitrary, 
    $f$ is continuous.
\end{proof}

Suppose $\tau_1$ and $\tau_2$ are two topologies on 
the same set $X$. 

We say $\tau_1$ is \textbf{finer} than $\tau_2$ if 
\[
    \tau_2 \subseteq \tau_1. 
\]
That means $\tau_1$ has more open sets. 
Equivalently, $\tau_2$ is \textbf{coarser} than 
$\tau_1$. For every set $X$: 
\[
    \{\emptyset, X\} \subseteq \tau \subseteq \mathcal P(X). 
\]
So the indiscrete topology is the coarsest topology possible, 
and the discrete topology is the finest topology possible. 

This affects continuity. If the domain topology is finer, 
continuity is easier. If the target topology is finer, continuity 
is harder. 

\begin{ex}
    Let $X$ have the discrete topology and let $Y$ be any 
    topological space. 
    Then every $f : X \to Y$ is continuous. 

    This is because $f^{-1}(V) \subseteq X$ and 
    every subset of $X$ is open. 
\end{ex}
\begin{ex}
    Let $Y$ have the indiscrete topology. Then 
    every function $f :X \to Y$ is continuous. 
    
    The preimages 
    \[
        f^{-1}(\emptyset) = \emptyset, \ 
        f^{-1}(Y) = x
    \]
    are both open in $X$. 
\end{ex}
\begin{ex}
    Let $X$ be a set with two topologies $\tau_1, \tau_2$. 
    Then $\operatorname{id} : (X, \tau_1) \to (X,\tau_2)$ is 
    continuous exactly when $\tau_2 \subseteq \tau_1$. 
\end{ex}

\subsection*{First glimpse of homeomorphisms}

Topology is about spaces up to continuous deformation. 
The correct notion of "same space" is not equality, but 
\textbf{homeomorphism}. 

\begin{defn}[Homeomorphism] 
    Let $X,Y$ be topological spaces. 

    A function $X \to Y$ is a \textbf{homeomorphism} if 
    \begin{enumerate}
        \item $f$ is bijective, 
        \item $f$ is continuous, 
        \item $f^{-1}$ is continuous.
    \end{enumerate}
    If such an $f$ exists, then $X$ and $Y$ are homeomorphic. 
\end{defn}
Intuitively, homeomorphic spaces have the same topological shape. 
Think about those videos of coffee cups and donuts. This is what 
they meant. 

\begin{ex}
    $(0,1)$ and $\mathbb R$ are homeomorphic. 
    One possible homeomorphism is 
    \[
        f : (0,1) \to \mathbb R, \ 
        x \overset{f}{\mapsto} \tan\left(\pi x - \frac{\pi}{2}\right).
    \]
\end{ex}
So $\mathbb R$ is homeomorphic to $(0,1)$ even though $(0,1)$ is 
"small" (bounded) in comparison to $\mathbb R$.

Lastly we remark that openness doesn't need to have 
anything to do with distance. 
In algebraic geometry and commutative algebra we look at 
the prime spectrum of a ring, 
\[
    \operatorname{Spec}(A)
\]
which we endow with a topology called the \textbf{Zariski topology}. 
Closed sets are defined by ideals: 
\[
    V(I) = \{\mathfrak p \in \operatorname{Spec}(A) : 
    I \subseteq \mathfrak p\}. 
\]
So the closed subsets are controlled by equations and ideals. 

\subsection*{Exercises}

\begin{prob}
    Let $X$ be any set. Prove that 
    the discrete topology $\tau = \mathcal P(X)$ is 
    a topology on $X$. 
\end{prob}
\begin{proof}
    $\tau \subseteq \mathcal P(X)$ is ceirtainly satisfied. 
    Furthermore $\emptyset, X \in \tau = \mathcal P(X)$ is 
    also satisfied. 

    Now let $I$ be some index set. Since 
    $\mathcal P(X)$ contains all subsets of $X$ and we 
    necessarily have 
    \[
        \bigcup_{i \in I} U_i \subseteq X
    \]
    for $U_i \in \mathcal P(X)$, and then 
    we have 
    \[
        \bigcup_{i\in I} U_i \in \mathcal P(X) = \tau. 
    \]
    Thus $\tau$ is closed under arbitrary intersections. 

    Let $U_1,\ldots,U_n \in \tau = \mathcal P(X)$. 
    Once again since $U_i \subseteq X$ for all $1 \leq i \leq n$ we 
    have 
    \[
        \bigcap_{i\in[n]} U_i \subseteq X
    \]
    so 
    \[
        \bigcap_{i\in[n]} U_i \in \tau = \mathcal P(X).
    \]
\end{proof}

\begin{prob}
    Let $X$ be any set. Prove that $\tau = \{\emptyset, X\}$ 
    is a topology on $X$. 
\end{prob}
\begin{proof}
    We have $\emptyset, X \in \tau$ by definition and also 
    $\tau \subseteq \mathcal P(X)$. 

    Now notice that the only non-trivial collection 
    of open sets in $\tau$ is the entireity of 
    $\tau$: 
    \[
        \emptyset, X,
    \]
    and this collection is finite. 
    Then $\emptyset \cup X = X \in \tau$ and 
    $\emptyset \cap X = \emptyset \in \tau$. 

    Thus $\tau$ is closed under arbitrary unions 
    and finite intersection adn hence forms a topology 
    on $X$. 
\end{proof}
\begin{prob}
    Let $X$ be a topological space. Prove that arbitrary intersections 
    of closed sets are closed. 
\end{prob}
\begin{proof}
    Let $I$ be some index set so that 
    $C_i$ is closed for every $i \in I$. 

    Then $X \setminus C_i$ is open for every $i \in I$. 
    Then 
    \[
        \bigcup_{i\in I}X \setminus C_i
        = \bigcup_{i\in I} X \cap C_i^c 
    \]
    is open. 
    Taking the complement we get 
    \[
        \left(\bigcup_{i\in I} X \cap C_i^c\right)^c
        = \bigcap_{i\in I} \emptyset \cup C_i 
        = \bigcap_{i \in I} C_i
    \]
    so 
    \[
        \bigcap_{i\in I} C_i  
    \]
    is closed. 
\end{proof}

\begin{prob}
    Let $\tau_1,\tau_2$ be topologies on $X$. Prove that 
    $\operatorname{id} : (X, \tau_1) \to (X, \tau_2)$ is 
    continuous if and only if $\tau_2 \subseteq \tau_1$. 
\end{prob}
\begin{proof}
    Suppose $\tau_2 \subseteq \tau_1$. 
    Let $U \in \tau_2$. Then $U \in \tau_1$ and 
    \[
        \operatorname{id}^{-1}(U) = U \in \tau_1.
    \]
    Thus the preimage of $U \in \tau_2$ is in $\tau_1$. 
    As $U$ was arbitrary we have that $\operatorname{id}$ 
    is continuous. 

    Suppose $\operatorname{id}$ is continuous. Then for any open set 
    $U \in \tau_2$ we have that 
    \[
        \operatorname{id}^{-1}(U) = U \in \tau_1
    \]
    hence 
    \[
        \tau_2 \subseteq \tau_1.
    \]
\end{proof}

\subsection{Some exercises from Section 13 of Munkres}
\begin{prob}[Problem 1]
    Let $X$ be a topological space; let $A$ be a subspace of 
    $X$. Suppose that for each $x \in A$ there is an open set 
    $U$ containing $x$ such that $U \subset A$. Show that 
    $A$ is open in $X$. 
\end{prob}
\begin{proof}
    For each $x\in A$, let $U_x$ denote the open 
    set containing $x$ such that $U_x \subset A$. 
    We claim that 
    \[
        A = \bigcup_{x \in A} U_x.
    \]
    For each $x \in A$, $U_x \subset A$ so 
    \[
        \bigcup_{x \in A} U_x \subseteq A.
    \]
    Now, for some arbitrary $x \in A$, 
    there exists an open set $U_x$ such that 
    $x \in U_x$. Thus 
    \[ 
        A \subseteq \bigcup_{x \in A} U_x.
    \] 
    The two inclusions give the result.
\end{proof}

\begin{prob}[Problem 3]
    Show that the collection $\mathcal T_c$ given in 
    Example 4 of \S{}12 (of Munkres) is a topology 
    on the set $X$. 

    Is the collection 
    \[
        \mathcal T_\infty = \{U \bigm| X \setminus U \ 
        \text{is infinite or empty or all of} \ X\}
    \]
    a topology on $X$?

    Note: For a set $X$, define 
    $\mathcal T_c$ to be the collection of all 
    subsets $U$ of $X$ such that $X \setminus U$ is 
    either countable or all of $X$. 
\end{prob} 
\begin{proof}
    Since $X \setminus \emptyset = X$ so 
    $\emptyset \in \mathcal T_c$. 
    Furthermore, since $X - X = \emptyset$ which is 
    finite and thus countable, we have 
    $X \in \mathcal T_c$. 

    Let $\{U_\alpha\}$ be a family of open sets 
    (ignoring the trivial case where $U_\alpha = \emptyset$ for 
    every such $U_\alpha$)
    in $\mathcal T_c$. By definition, 
    for each $U$, we have that 
    \[
        |X \setminus U_\alpha| \leq \aleph_0. 
    \]
    Since 
    \[
        \left|\bigcup_{\alpha} U_\alpha\right| \geq |U_\alpha|
    \]
    for any $\alpha$, we have 
    \[
        \left|X \setminus \bigcup_{\alpha}U_\alpha\right| 
        \leq |X \setminus U_\alpha| \leq \aleph_0. 
    \]
    Hence $\bigcup_{\alpha} U_\alpha \in \mathcal T_c$ so 
    $\mathcal T_c$ is closed under arbitrary unions. 

    Now let $\{U_i\}_{i\in[n]}$ be a finite family of open sets. 
    Assume also that $U_i \neq \emptyset$ for every $i$ as that 
    case is trivial. Then 
    \[
        X \setminus \bigcap_{i=1}^n U_i 
        = \bigcup_{i=1}^n (X\setminus U_i). 
    \]
    By assumption, $X\setminus U_i$ is countable and since the 
    union of countable sets is countable we have 
    that $\mathcal T_c$ is closed under finite intersections. 

    Now to answer the second question, $\mathcal T_\infty$ is 
    absolutely not a topology on $X$. Take for instance 
    $X = \mathbb N$. Then $\{n\} \in \mathcal T_\infty$, 
    but 
    \[
        \mathbb N \setminus \bigcup_{n \in \mathbb N \setminus \{7\}} \{n\}  
    \]
    is $\{7\}$ which is neither infinite, empty, or all of $X$. 
\end{proof}

\begin{prob}[Problem 4]
    \medskip 
    \
    \begin{enumerate}[label=\alph*)]
        \item If $\{\mathcal T_\alpha\}$ is a 
            family of topologies on $X$, show that $\bigcup \mathcal T_\alpha$
            is a topology on $X$. Is $\bigcup \mathcal T_\alpha$ 
            a topology on $X$? 
    \end{enumerate}
\end{prob}
\begin{proof}[Proof of (a)]
    Suppose $\{\mathcal T_\alpha\}$ is a family of 
    topologies on $X$. Then 
    $\emptyset, X \in \mathcal T_\alpha$ for all $\alpha$ so 
    \[
        \emptyset, X \in \bigcap \mathcal T_\alpha. 
    \]

    Let $\{U_\beta\}$ be a family of open sets in 
    $\bigcap \mathcal T_\alpha$. Then 
    every $U_\beta$ lies in $\mathcal T_\alpha$ for every $\alpha$. 
    Thus 
    \[
        \bigcup U_\beta \in \mathcal T_\alpha, \ \forall \alpha
    \]
    and hence $\bigcup U_\beta$ lies in the intersection. 

    Now let $\{U_i\}_{i\in[n]}$ be a finite family of open sets in 
    $\bigcap \mathcal T_\alpha$. 

    Again, $U_i \in \mathcal T_\alpha$ for every $i$ and $\alpha$ and 
    the argument is analogous to the one above. 

    The same does not hold for $\bigcup \mathcal T_\alpha$. 
    Suppose $X = \{1,2,3\}$ with 
    $\mathcal T_1 = \{\emptyset, X, 1\}$ and $\mathcal T_2 = 
    \{\emptyset, X, 2\}$. Then 
    $\{1\}$ and $\{2\}$ lie in their union, but $\{1,2\}$ does not. 
\end{proof}

\begin{prob}[Problem 6]
    Show that the topologies $\mathbb R_\ell$ and $\mathbb R_K$ 
    are not comparable. 
\end{prob}
\begin{proof}
    Recall that the topology of $\mathbb R_\ell$ is 
    \[
        \mathcal B_\ell = \{[a,b) \bigm| a < b\},
    \]
    and that the topology of $\mathbb R_K$ is 
    \[
        \mathcal B_K = \{(a,b) : a < b\} \cup \{(a,b) \setminus K 
        : a < b\} 
    \]
    where 
    \[
        K = \left\{\frac{1}{n} \biggm| n \in \mathbb Z_+\right\}.
    \]
    These are not comparable. Take for instance 
    $(-1,1) \setminus K \in \mathcal B_K$. If $\mathcal B_K \leq 
    \mathcal B_\ell$ we should be able to find some $\varepsilon > 0$ so that 
    $[0,\varepsilon) \subseteq (-1,1) \setminus K$, but 
    for every $\varepsilon > 0$ we can choose $n$ sufficiently large 
    so that 
    \[
        \frac{1}{n} < \varepsilon. 
    \]
    Then $1/n \in [0,\varepsilon)$ implying $[0,\varepsilon) \not\subseteq 
    (-1,1) \setminus K$. Thus $\mathcal B_K \not\subseteq \mathcal B_\ell$ 

    For the other way consider $[0,1) \in \mathcal B_\ell$. 
    Every $K$-neighbourhood at $0$ is either $(a,b)$ or $(a,b) \setminus K$ 
    with $a < 0 < b$. Either way it contains some negative real numbers so 
    no such basic neighborhood of $0$ can fit in $[0,1)$. 
\end{proof}

\begin{prob}[Problem 8]
    \medskip 
    \ 
    \begin{enumerate}[label=\alph*)]
        \item Apply Lemma 13.2 to show that the 
            countable collection 
            \[
                \mathcal B = \{(a,b) \bigm| a < b, \ a,b \in \mathbb Q\}
            \]
            is a basis that generates the standard topology on 
            $\mathbb R$.
    \end{enumerate}
\end{prob}
\begin{proof}[Proof of (a)]
    Recall that Lemma 13.2 says that if $X$ is a topological space, and 
    $\mathcal C$ is a collection of open sets such that for each 
    open set $U$ of $X$ and each $x \in U$ there is some element  
    $C \in \mathcal C$ such that $x \in C \subset U$, then $\mathcal C$ 
    is a basis for the topology on $X$.

    Let $U$ be open in the canonical topology on $\mathbb R$. Then 
    $U = (a_0,b_0)$ for some $a_0,b_0 \in \mathbb R$ such that 
    $a_0 < b_0$. By virtue of $\mathbb Q$ being dense in $\mathbb R$
    we can, for each $\varepsilon > 0$ find some $a,b \in \mathbb Q$ 
    suitably close: 
    \[
        a - a_0 < \varepsilon, \ b_0 - b < \varepsilon. 
    \]
    For any $x \in U = (a,b)$ let $\varepsilon = \min\{x - a_0, b_0-x\} > 0$. 
    Choose rational $a$ and $b$ (with $a < b$) so that 
    $a - a_0, b_0 - b < \varepsilon$. 
    Then 
    \begin{align*}
        a - a_0 &< \varepsilon \\ 
                &\leq x - a_0 \\ 
        a       &< x, 
    \end{align*}
    and 
    \begin{align*}
        b_0 - b &< \varepsilon \\
                &\leq b_0 - x \\ 
        -b      &\leq -x \\ 
        b       &> x. 
    \end{align*}
    Thus, by Lemma 13.2, $\mathcal B$ is a basis for the standard 
    $\mathbb R$ topology. 

    This is a bit imprecise in places, but the idea is right.
\end{proof}

\section{Some problems to recap}

\begin{prob}[Munkres \S 16, Exercise 3]
    Let 
    \[
        [-1,1] \subseteq \mathbb R
    \]
    with the subspace topology inherited from 
    the standard topology on the reals. 

    For each of the following sets, determine: 
    \begin{enumerate}
        \item Is it open in $Y$? 
        \item Is it open in $\mathbb R$? 
    \end{enumerate} 
    \begin{itemize}
        \item $A = \left\{x \biggm| \frac{1}{2} < |x| < 1\right\}$, 
        \item $B = \left\{x \biggm| \frac{1}{2} < |x| \leq 1\right\}$, 
        \item $C = \left\{x \biggm| \frac{1}{2} \leq |x| < 1\right\}$, 
        \item $D = \left\{x \biggm| \frac{1}{2} \leq |x| \leq 1\right\}$, 
    \end{itemize}
    and 
    \begin{itemize}
        \item $E = \{x \bigm| 0 < |x| < 1 \ \text{and} \ 1/x
        \not\in \mathbb Z_+\}$.
    \end{itemize}
\end{prob}
\begin{proof}[Sol]
    A set $X$ is open in $\mathbb R$ if for every 
    $x \in X$, there is some $\varepsilon > 0$ 
    such that $(x-\varepsilon, x + \varepsilon) \subseteq X$. 
    This is the same behavior inherited by $Y$. 

    We begin with $A$. 
    We claim 
    that this set is indeed open. 

    We have that 
    \[
        \frac{1}{2} < x < 1 \ \text{or} \ 
        -1 < x < -\frac{1}{2}.  
    \]
    Hence $A = (-1,-1/2) \cup (1/2, 1)$. The union of 
    open sets is open hence $A$ is open. 

    Now we do $B$. We claim that this set is not open. 
    Let $x = 1$. There would have to exist $\varepsilon > 0$
    so that $(x - \varepsilon, x + \varepsilon) \subseteq X$, but 
    $x + \varepsilon > 1$ and hence not in $B$. 

    Now we do $C$. We also claim this set is not open 
    using the same argument. 
    An open $\varepsilon$-ball of $1/2$ would include 
    $1-\varepsilon/2 < 1/2$. 

    $D$ fails for the same reason. 

    \[
        E = (-1,0) \cup \bigcup_{n=1}^\infty 
        \left(\frac{1}{n+1}, \frac{1}{n}\right),
    \]
    a union of open sets, hence $E$ is open.
\end{proof}

\begin{prob}
    Let 
    \[
        f : X \to Y
    \]
    and suppose $\mathcal B$ is the 
    basis for the topology on $Y$.

    Prove that $f$ is continuous if 
    and only if 
    \[
        f^{-1}(B)
    \]
    is open for every $B \in \mathcal B$.
\end{prob}
\begin{proof}
    Suppose $f$ is continuous. 
    Since $B \in \mathcal B$ are necessarily 
    open in the topology they induce it is trivial 
    that 
    \[
        f^{-1}(B)
    \]
    is open, by definition. 

    Now suppose 
    \[
        f^{-1}(B) = \{x \in X \bigm| f(x) \in B\}
    \]
    is open for every $B \in \mathcal B$. 
    Recall that every open set $U \subseteq Y$ can be 
    written as 
    \[
        U = \bigcup_{\alpha} B_\alpha, \ B_\alpha \in \mathcal B. 
    \]
    Then 
    \[
        f^{-1}\left(\bigcup_\alpha B_\alpha\right) 
        = \{x \in X \bigm| x \in f(B_\alpha) \ \text{for at least one} \ \alpha\}. 
    \]
    In particular 
    \[
        f^{-1}\left(\bigcup_\alpha B_\alpha\right)
        = \bigcup_\alpha f^{-1}\left(B_\alpha\right)
    \]
    and the preimage of $B_\alpha$ is, by assumption, 
    open in $X$. As topologies are closed under arbitrary 
    unions we get that $f^{-1}(U)$ is open in $X$. 
    As $U$ was arbitrary, this holds for any open $U \subseteq Y$. 

\end{proof}

\begin{prob}
    Let $X,Y,Z$ be topological spaces and let 
    \[
        f : Z \to X \times Y. 
    \]
    Write 
    \[
        f(z) = (f_1(z), f_2(z)) 
    \]
    where $f_1,f_2$ are the projections 
    to coordinates $1$ and $2$ respectively. 

    Prove that $f$ is continuous if and only if 
    $f_1,f_2$ are continuous.
\end{prob}
\begin{proof}
    The basis of the product topology is 
    \[
        \{U \times V \bigm| U \subseteq X \ 
        \text{is open, and} \ V \subseteq Y \ \text{is open}\}.
    \]

    Suppose $f$ is continuous. Then 
    \[
        f_1 = \pi_1 \circ f, \ \text{and} \ 
        f_2 = \pi_2 \circ f
    \]
    where $\pi_1,\pi_2$ are projections. 
    Thus $f_1,f_2$ are compositions of continuous functions 
    which is continuous. 

    Now suppose $f_1,f_2$ are continuous. 

    An open set $W \in Y$ can be written as 
    \[
        W = \bigcup_{\alpha} U_\alpha \times V_\alpha
    \]

    Then 
    $f^{-1}(W)$ is the set of all $z \in Z$ 
    such that $f_1(z) \in U_\alpha$ for at least one 
    $\alpha$ and $f_2(z) \in V_\alpha$ for at least one 
    alpha so 
    \[
        f^{-1}(W) = \bigcup_\alpha f_1^{-1}(U_\alpha) \cap f^{-1}_2(V_\alpha). 
    \]
    As the $f_1^{-1}(U_\alpha), f_2^{-1}(V_\alpha)$ are open 
    in $Z$, their finite union is open in $Z$ and 
    hence the union of $f_1^{-1}(U_\alpha), 
    f_2^{-1}(V_\alpha)$ over $\alpha$ is 
    open in $Z$.
\end{proof}

\begin{prob}
    Let $\mathcal T, \mathcal T'$ be topologies on a set 
    $X$ with $\mathcal T \subseteq \mathcal T'$, meaning 
    $\mathcal T$ is finer than $\mathcal T'$. 
    Considering the identity map 
    \[
        \operatorname{id}: X \to X
    \]
    we can regard it in two ways, the one 
    where $\mathcal T'$ is carried to $\mathcal T$ 
    or the other way. 
    For which of these is the identity map always 
    continuous?
\end{prob}
\begin{proof}
    The map 
    \[
        \operatorname{id} : (X, \mathcal T') \to 
        (X, \mathcal T)
    \]
    is always continuous, the other doesn't 
    need to be. 

    Let $U \in \mathcal T$ be open. Then 
    \[
        \operatorname{id}^{-1}(U) = U \in \mathcal T'. 
    \]

    Consider $\mathcal T = \{\emptyset, X\} \subseteq 
    \mathcal T' = \{\emptyset, X, A\}$. 
    Then the other identity map has 
    \[
        \operatorname{id}(A)^{-1} = A \not\in \mathcal T. 
    \]
\end{proof}

\section{Quotient spaces and quotient topology}

Quotient topology will be in spirit antithetical to 
the product topology. The fundamental question 
is: if we glue points of some space $X$ together, what 
topology should the space have? 

We begin by recapping what equivalence relations are. 

\begin{defn}
    Let $X$ be a set. 

    An equivalence relation, 
    on $X$ is relation $\sim$ satisfying reflexivity: 
    \[
        x \sim x
    \]
    for every $x\in X$, 
    
    symmetry: 
    \[
        x \sim y \Leftrightarrow y \sim x, 
    \]
    and 

    transitivity: 
    \[
        x \sim y, \ y \sim z \Rightarrow x \sim z. 
    \]

    For $x \in X$ we denote the equivalence class 
    \[
        [x] = \{y \in X \bigm| y \sim x\}. 
    \]
\end{defn}

Equivalence classes partition $X$. 

The quotient set is 
\[
    X/\sim \ = \{[x] \bigm| x \in X\}. 
\]

There is a canonical map 
\[
    q : X \to X/\sim \ , \ q(x) = [x].
\]
This map is automatically surjective. 

This is a nice construction in topology because 
these relations give us a sort of 
language for gluing. 

Take $X = [0,1]$. Suppose we identify two endpoints 
$0\sim 1$, while every other endpoint is 
equivalent to only itself. Then the quotient 
\[
    [0,1]/\sim
\]
is like a circle. 

Likewise starting from a square 
\[
    [0,1]\times [0,1]
\]
and identifying the left and right edges 
appropriately gives a cylinder. 

Now the natural question is what topology should be endowed 
on this space? 

Suppose $X$ is a topological space and let 
\[
    q : X \to Y
\]
be a surjection. 

We want to make $Y$ a topological space. 

Since $q$ tells us which points of $X$ collapse onto 
the same point in $Y$, the natural idea is 
a set downstairs should be open exactly when its 
full preimage upstairs is open. 

\begin{defn}[Quotient topology]
    The quotient topology on $Y$ is 
    \[
        \tau_Y = \{U \subseteq Y \bigm| q^{-1}(U) \ \text{open 
        in} \ X\}. 
    \]
\end{defn} 


\begin{prop}
    The quotient topology is indeed a topology. 
\end{prop}
\begin{proof}
    First, 
    \[
        q^{-1}(\emptyset), \ q^{-1}(Y) = X. 
    \]
    Both are open in $X$. 

    Let $\{U_\alpha\}_\alpha \subseteq \tau_Y$. Then 
    \[
        q^{-1}\left(\bigcup_\alpha U_\alpha\right) 
        = \bigcup_\alpha q^{-1}(U_\alpha). 
    \]
    Each $q^{-1}(U\alpha)$ is open in $X$ so 
    the union of these is open in $X$. 

    If $U_1,\ldots,U_n \in \tau_Y$, than 
    \[
        q^{-1}\left(\bigcup_{i=1}^n U_i\right)
        = \bigcap_{i=1}^n q^{-1}(U_i)
    \]
    which is the finite intersection of open sets in 
    the domain and thus open. 
\end{proof}

We need the map to be surjective otherwise 
there could be points in the codomain which 
are never seen. 

A surjective continuous map $q : X \to Y$ is called 
a quotient map if $U \subseteq Y$ if and only if $q^{-1}(U)$
is open in $X$. 

A quotient space, on the other hand, is 
typically a space 
\[
    X/\sim
\]
equipped with the quotient topology induced by 
the canonical projection $q : X \to X/\sim$. 

It is tempting to think that 
\[
    U \subseteq X \ \text{open} 
    \Rightarrow q(U) \subseteq Y \ \text{open}.
\]
This is false in general.

The definition works with the preimage not with the 
image. Just like how topology in general behaves quite 
nicely with preimages. 

\begin{defn}
    A subset $A \subseteq X$ is called saturated 
    with respect to $q$ if it contains 
    every fibre that it touches. 
    Formally, 
    \[
            A = q^{-1}(q(A)). 
    \]
    In the case of an equivalence relation this means 
    \[
        \text{If} \ x \in A, \ \text{then every 
        point is equivalent to} \ x \ \text{is also in} \ A.
    \]
\end{defn}

Every set of the form $q^{-1}(A)$ is saturated. 

Thus the quotient topology is really controlled by 
open saturated subsets of $X$. 

\begin{ex}
    Let $X$ be the closed unit interval. 
    Define $0 \sim 1$ with no other nontrivial identifications.

    Then $Y = X/\sim$ has one special point 
    \[
        [0] = [1]. 
    \]

    Suppose $U \subseteq Y$ contains $[0]$. 
    Its preimage must contain both $0$ and $1$. 
    For $U$to be open in $Y$ it must be open in 
    the subspace $[0,1]$. So it must contain 
    something like 
    \[
            [0,\varepsilon) \cup (1-\varepsilon, 1].
    \] 

    After quotienting, these two pieces become 
    one neighborhood around the glued point. 

    That is exactly what neighborhoods around a point 
    on the circle look like. 
\end{ex}

Now we arrive at the universal property. 

Suppose 
\[
    q : X \to Y
\]
is a quotient map. 

Let $Z$ be another topological space and 
\[
    f : X \to Y. 
\]

Then 
\[
    f \ \text{is continuous} \Leftrightarrow 
    f \circ g \ \text{is continuous}.
\]

You can illustrate the difference between 
this- and the product topology's universal property 
particularly nice with diagrams in a category theory-esque 
way, but I'm lazy. 

Now we pose a natural question. 

Suppose 
\[
    q : X \to X/\sim
\]
is the quotient map. 
Imagine we have a continuous function 
\[
    g : X \to Z.
\]

When can it descend to a function $\overline g: X/\sim \to Z$
such that $g = \overline g \circ q$? 

Set-theoretically, this is possible exactly when 
$g$ is constant on equivalence classes: 
\[
    x \sim y \Rightarrow g(x) = g(y). 
\]

If that happens, define 
\[
    \overline g([x]) = g(x). 
\]

Because $g$ is constant on equivalence classes, 
this is well-defined. 

Then the quotient universal property gives 
\[
    g = \overline g \circ q
\]
continuous 
\[
    \Rightarrow \overline g \ \text{continuous}. 
\]

So quotient spaces let you convert functions on $X$ 
that respect the identification into functions 
on the glued space. 

This is perhaps the most useful operational idea: 

\begin{align*}
    \text{If a continuous map cannot distinguish points 
    you glue together,} \\  
    \text{it descends continuously 
    to the quotient}. 
\end{align*}

\begin{prob}
    Let $X = [0,1]$ and define $0 \sim 1$ with 
    all other equivalence classes singletons. 
    Let 
    \[
        q : X \to X/\sim
    \]
    be the quotient map. 
    Consider 
    \[
        A = q([0,1/4) \cup (3/4,1]). 
    \]
    Try to determine whether $A$ is open 
    in the quotient space. 
\end{prob}

\begin{proof}[Sol.]
    Let $B = [0,1/4) \cup (3/4, 1]$. The forward inclusion 
    \[
        B \subseteq q^{-1}(q(B)) 
    \]
    is always true. Indeed, if $x \in B$ then $q(x) \in q(B)$ 
    so 
    \[
        x \in q^{-1}(q(B)).
    \]

    For the reverse, suppose $x \in q^{-1}(q(B))$. 
    Then 
    \[
        q(x) \in q(B), 
    \]
    so there exists $b \in B$ such that 
    \[
        q(x) = q(b). 
    \]
    But $q(x) = q(b)$ means that 
    \[
        x \sim b. 
    \]
    The only nontrivial class is $\{0,1\}$. 
    Since $B$contains both $0,1$ whenever 
    $b \in B$ every point is equivalent to $b$ in $B$ 
    so $x \in B$. 
    
    The two inclusions
    mean that $B$ is staurated under the equivalence 
    relation. Hence $A$ is open in the 
    quotient space. 
\end{proof}

\begin{prob}
    \medskip 
    \ 

    \begin{enumerate}[label=\alph*)]
        \item Let $p : X \to Y$ be a continuous map. 
            Show that if there is a continuous map 
            $f : Y \to X$ such that $p \circ f$ equals
            the identity map of $Y$, then $p$ is a quotient map. 
        \item If $A \subset X$, a retraction of 
            $X$ onto $A$ is a continuous map $r : X \to A$
            such that $r(a) = a$ for each $a \in A$. Show 
            that a retraction is a quotient map. 
    \end{enumerate}
\end{prob}
\begin{proof}[Proof of (a)]
    Recall that a quotient map $g : X \to Y$ is a 
    surjection such that $U \subseteq Y$ 
    is open if and only if $f^{-1}(U)$ is open. 

    Condition two is partially satisfied 
    by $p$ since, by assumption, 
    $p^{-1}(U)$ is open when $U \subseteq Y$ is open. 

    Suppose $p$ is not a surjection. 
    Then there are $y \in Y$ such that, 
    for every $x \in X$, 
    \[
        p(x) \neq y. 
    \]
    Then, $p \circ f \neq \operatorname{id}_Y$. 
    $\operatorname{im} f \subseteq X$ so there is 
    no $x \in \operatorname{im} f$ so that
    \[
        p(x) = y. 
    \]
    In particular, 
    $f(y) \in \operatorname{im} f \subseteq X$ so 
    $f(y) \neq y$. 

    Thus we require that $p$ is a surjection. 

    Now it remains to show that $p$ carries 
    open sets to open sets. 
    Suppose $V \subseteq X$ is open. 
    Then since $f$ is continuous, 
    $f^{-1}(V) \subseteq Y$ is open. 

    We have that 
    \[
        p\circ f(f^{-1}(V)) = f^{-1}(V)
    \]
    so 
    \[
        p(V) = f^{-1}(V). 
    \]
    Thus $p$ carries open sets to open sets. 

    Summing up, we have that $p$ is 
    a quotient map. 
\end{proof}
\begin{proof}[Proof of (b)]

Let $r: X \to A$ be a retraction, and let
\[
    i\hookrightarrow X
\]
denote the inclusion map, defined by $i(a)=a$.

Since $A$ has the subspace topology, $i$ is continuous. Moreover, since $r$ is a retraction,
\[
    (r\circ i)(a)=r(a)=a
\]
for every $a\in A$. Hence
\[
    r\circ i=\operatorname{id}_A.
\]

Thus $r$ has a continuous right inverse. By part $a$, $r$ is a quotient map.
\end{proof}

\begin{prob}
    \medskip 
    \ 
    \begin{enumerate}[label=\alph*)]
        \item Define an equivalence relation on the plane 
            $X = \mathbb R^2$ as follows: 
            \[
                x_0 \times y_0 \sim x_1 \times y_1 \ \text{if} \ 
                x_0 + y_0^2 = x_1 + y_1^2. 
            \]
            Let $X^*$ be the corresponding quotient space. It is homeomorphic 
            to a familiar space; what is it? 
        \item Repeat for the equivalence relation 
            \[
                (x_0, y_0) \sim (x_1,y_1) \ \text{if} \ 
                x_0^2 + y_0^2 = x_1^2 + y_1^2. 
            \]
    \end{enumerate}
\end{prob}
\begin{proof}[Solution to (a)]
    Define $g([x,y]): X^* \to \mathbb R, \ g([x,y]) = x+y^2$. 

    Then, 
    \[
        (x_0,y_0) \sim (x_1,y_1) 
        \Leftrightarrow 
        g([x_0,y_0]) = g([x_1,y_1]). 
    \]
    Thus the equivalence classes are 
    the fibres of $g$.
    
    We claim this map is the continuous 
    inverse of the quotient map 
    $q : \mathbb R^2 \to \mathbb R$. 

    Well-definedness of the map is 
    immediate. 

    It is also injective immediately from the definition. 

    For surjectivity, given $t \in \mathbb R$, simply take 
    $(t,0)$. Then $g([t,0]) = t$. 

    Notice $x + y^2 = g \circ q$ is continuous. 
    By the definining property of the quotient topology, 
    $g$ is continuous exactly when $g \circ q$ is. 
    Thus $g$ is continuous. 

    Define $s : \mathbb R \to \mathbb R^2$ 
    by $s(t) = (t,0)$. Then $s$ is continuous. 

    Now $(q \circ s) (t) = [(t,0)]$, but this is 
    exactly the inverse of $g$.
    So $g^{-1}$ is continuous. 

    Thus $X^* \cong \mathbb R$. 
\end{proof}
\begin{proof}[Solution sketch of (b)]
    Proceed as above and define 
    $g : \mathbb R^2 \to [0,\infty)$ by 
    $g(x,y) = x^2 + y^2$ to 
    get $X^* \cong [0,\infty)$.
\end{proof}

\subsection{Some more problems}

\begin{prob}
    Let $q: X \to Y$ be a surjection. 
    Prove that for every $A \subseteq X$, 
    \[
        A \subseteq q^{-1}(q(A)). 
    \]
    Show that equality holds if and only if 
    $A$ is a union of fibers of $q$. 
    In the equivalence-relation setting, formulate 
    this as ''$A$ is a union of equivalence classes''. 
\end{prob}
\begin{proof}
    Let $A \subseteq X$. 
    Then $q(A) \subseteq Y$. 
    We look at the preimage 
    \[
        q^{-1}(q(A)) = \{x \in X \bigm| q(x) \in q(A)\}. 
    \]
    Clearly, this must include $A$. Otherwise 
    there would be at least one $a \in A$ such that 
    $q(a)\not\in q(A)$, but $q(A) = \{q(a) \bigm| a \in A\}$ 
    (Arguing about sets is ridiculous).

    Suppose the preimage is exactly $A$, meaning 
    $A$ is saturated w.r.t $q$. 
    Take any $a \in A$. Consider the fiber of $q$ 
    through $a$: 
    \[
        q^{-1}(\{q(a)\}). 
    \]
    We claim this entire fiber lies in $A$. 
    
    Indeed, if 
    \[
        x \in 
        q^{-1}(\{q(a)\}), 
    \]
    then 
    \[
        q(x) = q(a). 
    \]
    Since $a \in A$, 
    \[
        q(a) \in q(A),  
    \]
    so $q(x) \in q(A)$. 

    Thus 
    \[
        x \in q^{-1}(q(A)). 
    \]

    By assumption, $q^{-1}(q(A)) = A$, so $x \in A$. 

    Conversely, suppose $A$is a union of 
    fibers. 
    Then it remains to prove that 
    \[
        q^{-1}(q(A)) \subseteq A. 
    \]
    Take $x \in q^{-1}(q(A))$. 
    Then $q(x) \in q(A)$, 
    so by definition of the image there is some 
    $a \in A$ such that $q(x) = q(a)$. 
    Thus $x$ and $a$ lie in the same fiber. 
    Hence 
    \[
        q^{-1}(q(A)) \subseteq A. 
    \]
\end{proof}

\begin{prob}
    Let $X \subseteq \mathbb R$ be the 
    closed unit interval with $0 \sim 1$ 
    and define $g : [0,1] \to \mathbb R^2$ as 
    \[
        g(t) = (\cos 2\pi t, \sin 2\pi t). 
    \]
    Show that $g$ is constant on equivalence classes, 
    so it induces a well-defined map 
    \[
        \overline g : X / \sim \to S^1.
    \]
    Then use the quotient universal property to 
    prove $\overline g$ is continuous. 
\end{prob}

\begin{proof}[Proof sketch]
    $\cos$ and $\sin$ are $2\pi$-periodic so that 
    \[
        \cos(2\pi \cdot 0) = \cos(2\pi \cdot 1) \ \text{and} 
        \ \sin(2\pi \cdot 0) = \sin(2\pi \cdot 1). 
    \]
    This, together with the rest of the equivalence classes 
    being singletons, makes $g$ constant on 
    all the equivalence classes. 
    Thus $\overline g : X / \sim \ \to S^1$ is well-defined. 
    
    As this map is the composition of 
    $g$ (continuous) 
    with the canonical projection $x \overset\pi\mapsto 
    [x]$ then, by the universal property of quotient topology, 
    $\overline g$ is continuous. 

    Furthermore $\overline g$ is invertible (I will not 
    rigorously prove this) so 
    \[
        X/ \sim \ \cong S_1
    \]
    the $1$-dimensional loop. 
\end{proof}

\begin{prob}
    Let $X = \mathbb R$ and identify 
    all integers to a single point, leaving all 
    nonintegers equivalent only to themselves. 
    Give $X / \sim$ the quotient topology. 
    Decide whether the quotient is Hausdorff. 
\end{prob}
\begin{proof}
    For the quotient map $q : X \to X/\sim$ we 
    have $A \subseteq X/\sim$ is open if and only 
    its preimage is open in $X$. 

    The Hausdorff property is that 
    for any two distinct points $x,y$ we can 
    find open neighborhoods for each which are disjoint.

    We split it into two cases. 

    Case 1: $[\mathbb Z]$ and a noninteger point 

    Let $x \not\in \mathbb Z$. Since $\mathbb Z$ is 
    closed and discrete in $\mathbb R$, we can choose 
    $\varepsilon > 0$ such that 
    \[
        [x - \varepsilon, x +\varepsilon] \cap 
        \mathbb Z = \emptyset. 
    \]
    Put 
    \[
        V = (x - \varepsilon/2, x + \varepsilon / 2)
    \]
    and
    \[
        U = R / [x - \varepsilon, x + \varepsilon].
    \]
    Then $U$ and $V$ are disjoint w.r.t $q$. 
    Therefore 
    \[
        q^{-1}(q(U)) = U, \ 
        q^{-1}(q(V)) = V. 
    \]
    Since $U,V$are open, the quotient topology 
    implies that $Q(U), Q(V)$are open in 
    $\mathbb R/ \sim$. 
    They are also disjoint, because $U,V$ are disjoint 
    saturated sets. Finally 
    \[
        [\mathbb Z] \in q(U)
    \]
    and 
    \[
        [x] \in q(V). 
    \]
    Thus $[\mathbb Z], [x]$ have disjoint open neighborhoods. 

    The case for two nonintegers is easy from 
    the typical topology on $\mathbb R$. 
\end{proof}

\newpage 
\section{Connectedness, components, compactness}

From $\S 23$ of Munkres recall that a space 
$X$ is connected if we cannot write 
$X = U \sqcup V$ where $U,V$ are disjoint, nonempty, 
open subsets of $X$. 

Equivalently, 
$X$ is connected if and only if the only 
clopen subsets of $X$ are $X$ and $\emptyset$. 

Also, $X$ connected and $f : X \to Y$ continuous 
implies $f(X)$ connected. 

\subsection{\S 24, Connected subsets of the real line}
The main theorem we will be keeping in mind is 
$A \subseteq \mathbb R$ is connected if and only if 
$A$ is an interval, which is perhaps unsurprising.  

Interval should be understood in the broad order-theoretic 
sense, meaning 
\[
    x,z \in A, \ x < y < z \Rightarrow y \in A. 
\]
Such sets are also called convex subsets of 
the ordered set $\mathbb R$ in some contexts. 
Munkres develops this more generally for linear continua 
later, I believe. 

The claim now is that connected subsets must be intervals.

Suppose $A \subseteq \mathbb R$ is connected. 
Take $x,z \in A, \ x < z$. 
Suppose some $y$ satisfies $x < y < z$, but does 
not belong to $A$. 

Then 
\[
    A \cap (-\infty, y)
\]
and 
\[
    A \cap (y,\infty)
\]
are disjoint nonempty open subsets of $A$, and 
together they cover $A$. 

They are nonempty because $x$ is in the 
first and $z$ is in the second. 

Hence they give a separation of $A$, contradicting 
connectedness. 
Therefore every point between $x$ and $z$ lies in $A$. 
So $A$ is an interval. 

\subsection{The intermediate Value Theorem in topology}
Now suppose $X$ is connected and 
$f : X \to \mathbb R$ is continuous. 
Because continuous images of connected spaces 
are connected we have that $f(X)$ is connected, 
but connected subsets of $\mathbb R$ are intervals. 

Therefore if 
\[
    f(a) < r < f(b), 
\]
then 
\[
    r \in f(X). 
\]
So there exists $c \in X$ with $f(x) = r$. 

\subsection{Path connectedness}
Path connectedness is a strictly stronger 
notion that connectedness. 

A path from $x$ to $y$ in $X$ is a continuous 
map 
\[
    \gamma : [0,1] \to X
\]
such that 
\[
    \gamma(0) = x, \ \gamma(1) = y. 
\]
A space $X$ is path connected if every pair 
$x,y \in X$ can be joined by a path. 
Intuitively, connected says we cannot split the space, 
while path connected means we can travel 
continuously beetween all distinct points. 

\begin{prop}
    Path connectedness implies connectedness. 
\end{prop}
\begin{proof}
    Suppose $X$ is path connected, but disconnected. 
    Then $X = U \cup V$ is a separation. 
    Choose $x \in U$ and $y\in V$. 
    There is a path $\gamma : [0,1] \to X$ from 
    $x$ to $y$, but $[0,1]$ is connected, so its 
    continuous image $\gamma([0,1])$ is connected. 
    Yet this image intersects both $U$ and $V$, which 
    would separate it, a contradiction.  
\end{proof}

\subsection{\S 25, Components: how a disconnected space splits}
Connectedness is a yes/no question asking "Is the whole space 
connected?". 
Components are slightly more refined and instead ask 
"What are the largest connected chunks of the space?". 

For $x \in X$ define the components of $x$ to be the 
union of all connected subsets of $X$ containing $x$. 

Equivalently, define $x \sim y$ if there exists 
a connected subset of $X$ containing both $x$ and $y$. 
The equivalence classes are the components of $X$. So 
a component is a maximal connected subspace. 

Suppose 
\[
    C_x = \bigcup_\alpha A_\alpha, 
\]
where each $A_\alpha$ is connected and contains $x$. 
Every $A_\alpha$ has the common point $x$. 

The union of connected sets having a common point is 
connected. 

Therefore $C_x$ is connected. 
And since it contains every connected set 
containing $x$, it is maximal. 

Furthermore, components partition the space. 

Every point lies in its own component. Two components 
are either disjoint or equal. 
Suppose $C_1, C_2$ intersect. Then 
$C_1 \cup C_2$ is connected because the 
two connected sets have a common point. 
However, both $C_1,C_2$ are maximal connected sets so 
$C_1 = C_2$. 

Thus $X$ is partitioned into connected pieces. 

\subsection{Components are closed}
This is an important theorem: 
Every component of a space is closed. 

The main ingredient in this is an important one, 
namely that $A$ connected $\Rightarrow$ $\overline A$ connected. 
If $C$ is a component, then its closure $\overline C$ is connected, 
but $C \subseteq \overline C$ and $C$ was maximal so 
$\overline C = C$. Hence it is closed.  

\subsection{Path components} 
Similarly, define $x \sim_{p} y$ if there 
exists a path joining $x$ and $y$. 
The equivalence classes are called path components. 
Every path component lies inside a component because 
path connected $\Rightarrow$ connected. 

So conceptually, path components $\subseteq$ components and 
they can be strictly smaller. 

\subsection{Local connectedness} 

A space can faily to be connected globally, but look 
connected locally. 

A space $X$ is locally connected at $x$ if every 
neighborhood $U$ of $x$ contains a connected neighborhood 
$V$ of $x$. 

That is, 
\[
    x \in V \subseteq U
\]
with $V$ connected and a neighborhood of $x$. 

It is locally connected if this holds at every $x \in X$. 
Munkres gives an analogous definition for local path connectedness. 
For example, $\mathbb R \setminus \{0\}$ is not connected, 
but it is locally connected, meaning around every point 
you can choose a sufficiently small interval not crossing $0$. 

\subsection{Local theorem}

Munkres gives the useful characterization,  
\[
    X \ \text{is locally connected} 
\]
\[
    \Updownarrow 
\]
\[
    \text{every component of 
    every open subset of} \ X \ \text{is open in} \ X. 
\]

There is, of course, an analogous statement with 
path components and locally path connected spaces. 

Now, why should this make sense? 

Suppose $X$ is locally connected, $U \subseteq X$ is 
open, and $C$ is a component of $U$. Take $x\in C$. 

Because $U$ is a neighborhood of $x$, local connectednes 
gives a connected neighborhood $V_x$ with 
$x \in V_x \subseteq U$. 

Since $V_x$ is connected and intersects $C$, maximality 
of $C$ gives $V_x \subseteq C$. 
Do this for every $x \in C$. Then 
\[
    C = \bigcup_{x \in C} V_x. 
\]
Hence $C$ is open. 

Also, $X$ locally path connected $\Rightarrow$ components 
$=$ path components. 

\subsection{\S 26, Compactness}

Connectedness asks whether the space can be pulled 
apart, while compactness asks whether infinitely 
much local information can be reduced to finitely much 
information. 
Let $\mathcal U = \{U_\alpha\}_{\alpha \in A}$ be a collectioon 
of open sets in $X$. 

It is an open cover of $X$ if 
\[
    X \subseteq \bigcup_{\alpha \in A} U_{\alpha}. 
\]

The space $X$ is compact if every open cover has a 
finite subcover, meaning 
\[
    X \subseteq \bigcup_{\alpha \in A} U_{\alpha}
    \Rightarrow 
    X \subseteq U_{\alpha_1} \cup \cdots 
    \cup U_{\alpha_n}
\]
for some finite choice of indices. 

As a first example, recall from analysis that 
$[0,1]$ is compact, but $(0,1)$ is not. 
For example, $U_n = (1/n, 1)$ cannot be reduced to a 
finite cover. 

\subsection{Closed subsets of compact spaces}
\begin{thm}
    $X$ is compact and $A \subseteq X$ closed 
    implies that $A$ is compact. 
\end{thm}
\begin{proof}
    Let $\{U_\alpha\}$ be an open cover of $A$. 
    Because $A$ is closed, $X \setminus A$ is open. 
    Then $\{U_\alpha\} \cup \{X\setminus A\}$ covers $X$ 
    and can thus be reduced to a finite cover. 
    Discard $X \setminus A$ if it was chosen. The 
    remaining finitely many $U_\alpha$'s cover $A$. 
\end{proof}

\subsection{Compact subsets of Hausdorff spaces are closed} 
The converse is generally not true. 

In a Hausdorff space, $A \subseteq X$, $A$ compact, $X$ 
Hausdorff $\Rightarrow$ $A$ closed. 

Take $x \in X \setminus A$. For every $a \in A$, 
Hausdorffness gives disjoint open 
$U_a \ni a$ and $V_a \ni x$. 
The sets $U_a$ cover $A$. 
Compactness gives finitely many $U_{a_1}, \ldots, U_{a_n}$. 

Let 
\[
    V = V_{a_1} \cap \cdots \cap V_{a_n}. 
\]
This is open, contains $x$, and is disjoint 
from $U_{a_1} \cup \cdots \cup U_{a_n}$, which 
contains $A$. Thus $V \cap A = \emptyset$. So 
every point outside $A$ has an open neighborhood 
not intersecting with $A$. 
Therefore $X \setminus A$ is open, hence $A$ is closed. 

\subsection{Two nice theorems}
Another nice theorem I'll state informally: 
\[
    X \ \text{compact}, \ f: X \to Y \ \text{continuous} 
    \Rightarrow f(X) \ \text{compact. }
\]
Let $\{U_\alpha\}$ be an open cover of $f(X)$. 
Then $\{f^{-1}(U_\alpha)\}$ covers $X$. 
Compactness lets us reduce to finitely many and therefore 
we can find a finite cover of the image od $f$. 

\begin{thm}
    Suppose $f: X \to Y$ is continuous 
    and bijective with $X$ compact and $Y$ Hausdorff. 
    Then $f$ is a homeomorphism. 
\end{thm}
\begin{proof}[Proof sketch]
    Take a closed set $C \subseteq X$. Since 
    $X$ is compact, $C$ is compact. Continuous images 
    preserve that property so $f(C)$ is compact in $Y$. 
    Since $Y$ is Hausdorff, $f(C)$ is closed. Thus 
    $f$ maps closed sets to closed sets. 
    Since $f$ is a bijection, this implies $f^{-1}$ is continuous.
\end{proof}

\subsection{Some problems}
\begin{prob}[Munkres 24.1]
    \medskip
    \ 
    \begin{enumerate}[label=\alph*)]
        \item Show that no two of the spaces $(0,1), \ (0,1], 
            \ [0,1]$ are homeomorpnic. 

        \item Suppose there exist imbedddings $f : X \to Y$ and 
            $g : Y \to X$. Show by means of an 
            example that $X,Y$ need not be homeomorphic. 
        \item Show $\mathbb R^n$ and $\mathbb R$ 
            are not homeomorphic if $n > 1$. 
    \end{enumerate}
\end{prob}
\begin{proof}[Proof of (a)]
    Notice that $[0,1]$ is compact 
    while $(0,1)$ is not. 
    The latter is not since, if 
    we take for instance the cover 
    \[
        \{(1/n, 1)\}_{n\in \mathbb N}
    \]
    of $(0,1)$ this cannot be reduced to a 
    finite cover since for every $n$ we 
    can find some $m > n$ such that 
    \[
        0 < \underbrace{1/m}_{\in (0,1)} < 1/n.  
    \]
    So $[0,1]$ and $(0,1)$ cannot be homeomorphic since 
    continuous images would preserve compactness. 

    $[0,1]$ and $(0,1]$ cannot be homeomorphic for the same 
    reason. 

    Lastly $(0,1)$ is not homeomorphic to $(0,1]$. 
    Suppose they were. Let $f : (0,1] \to (0,1)$ be 
    the corresponding homeomorphism. 
    Then 
    \[
        f\bigm|_{(0,1] \setminus \{1\}} : 
        (0,1) \to (0,1) \setminus f(\{1\})  
    \]
    would be a homeomorphism. Let $x = f(1)$.  
    For any point in $y \in (0,1)$, its removal leaves  
    \[
        (0,1) \setminus \{y\} = (0,y) \cup (y,1)
    \]
    and in particular 
    \[
        (0,1) \setminus f(\{1\}) = (0,f(1)) \cup (f(1), 1)
    \]
    so the space is disconnected, while 
    $(0,1] \setminus \{1\} = (0,1)$ is connected, a 
    contradiction. 
\end{proof}

\begin{prob}[Munkres 25.4]
    Let $X$ be locally path connected. 
    Show that every connected open set in $X$ 
    is path connected.
\end{prob}

\begin{proof}
Let $U\subseteq X$ be open and connected. We show that
$U$ is path connected.

Fix $x\in U$ and let
\[
    P_x
    =
    \{z\in U : \text{there exists a path in $U$ from $x$ to $z$}\}.
\]
Thus $P_x$ is the path component of $x$ in $U$.

We first show that $P_x$ is open in $U$. Let $z\in P_x$.
Since $U$ is an open neighborhood of $z$ and $X$ is locally
path connected, there exists an open path-connected set
$V_z$ such that
\[
    z\in V_z\subseteq U.
\]
For every $w\in V_z$, there is a path from $z$ to $w$ inside
$V_z$. Since $z\in P_x$, there is also a path from $x$ to $z$
inside $U$. Composing these paths gives a path from $x$ to $w$,
so $w\in P_x$. Hence
\[
    V_z\subseteq P_x,
\]
and therefore $P_x$ is open.

We similarly show that $U\setminus P_x$ is open. Let
$z\in U\setminus P_x$. Again choose an open path-connected
neighborhood
\[
    z\in V_z\subseteq U.
\]
Suppose that $V_z\cap P_x\neq\emptyset$, and choose
$w\in V_z\cap P_x$. There is a path from $x$ to $w$, and since
$V_z$ is path connected, there is a path from $w$ to $z$.
Composing them would give a path from $x$ to $z$, contradicting
$z\notin P_x$. Hence
\[
    V_z\subseteq U\setminus P_x.
\]
Thus $U\setminus P_x$ is also open.

Therefore $P_x$ is both open and closed in $U$. Since $U$ is
connected and $P_x$ is nonempty, we must have
\[
    P_x=U.
\]
Hence every point of $U$ can be joined to $x$ by a path, so
$U$ is path connected.

\end{proof}

\begin{prob}
    Let $X$ be a topological space, and let 
    $A,B \subseteq X$ be connected subspaces such that 
    $A \cap B \neq \emptyset$. 
    Prove that $A \cup B$ is connected. 
\end{prob}
\begin{proof}
    We are free to discard the trivial case where 
    $A = B$. 

    Since $A,B$ are connected, 
    the only continuous functions $f : A \to \{0,1\}$ 
    and $g : B \to \{0,1\}$ are constant functions. 

    Suppose $A \cup B$ is not connected, meaning 
    we can find a continuous $h : A \cup B\to \{0,1\}$ which 
    is non-constant. 

    Necessarily, 
    \[
        h\bigm|_A = c_A \in \{0,1\}
    \]
    is constant, and 
    \[
        h\bigm|_B = c_B \in \{0,1\}
    \]
    is constant, since they are connected spaces. 

    We assume $h$ is non-constant. So there exists 
    some $x,y \in A \cup B$ such that 
    \[
        h(x) \neq h(y). 
    \]

    Then, necessarily, these points cannot be in 
    the intersection of $A$ and $B$. 
    So one of these points are in $A \setminus B$ or 
    $B \setminus A$. Without loss of generality, 
    suppose $x \in A \setminus B$. 
    Then we require $y \in B \setminus A$ since otherwise 
    we would have 
    \[
        c_A = h(x) \neq h(y) = c_A. 
    \]    
    From this we deduce $c_A \neq c_B$, but then, for 
    $z \in A \cap B$, we would have 
    \[
        h(z) = h\bigm|_A(z) = h\bigm|_B(z) 
    \]
    so 
    \[
        c_A = c_B, 
    \]
    a contradiction.
\end{proof}

\begin{prob}
    Let $\{A_\alpha\}_{\alpha \in I}$ be a 
    collection of connected subspaces of $X$. Suppose 
    there exists a point $p \in X$ such that 
    $p \in A_\alpha$ for every $\alpha \in I$. 

    Prove that 
    \[
        \bigcup_{\alpha \in I} A_\alpha 
    \]
    is connected. 
\end{prob}

\begin{proof}
    We suppose $A_\alpha$ are all connected, so 
    for any continuous $f_\alpha : A_\alpha \to \{0,1\}$, 
    the function must be constant. 

    Let $\varphi : \bigcup_{\alpha} A_\alpha \to \{0,1\}$ 
    be continuous. 

    Then, necessarily, 
    \[
        \varphi \bigm|_{A_\alpha} : A_\alpha \to     
        \{0,1\}
    \]
    is constant (also the restriction remains 
    continuous). 

    We have 
    \[
        \varphi(p) = c \in \{0,1\}
    \]
    and in particular, 
    since $p \in A_\alpha$ for every $\alpha \in I$, and 
    the restriction to $A_\alpha$ is constant, 
    we have 
    \[
        \im\varphi\bigm|_{A_\alpha} = \{c\} 
    \]
    for every $\alpha \in I$. Thus 
    \[
        \im\varphi = \{c\} 
    \]
    is constant, and hence 
    \[
        \bigcup_\alpha A_\alpha
    \]
    is connected.
\end{proof}

\begin{prob}
    Let $A\subseteq X$ be connected. Prove that its 
    closure $\overline A$ is connected. 
\end{prob}

\begin{proof}
    Suppose $A$ is a connected subspace of some 
    space $X$. Then any continuous function $f : A \to \{0,1\}$ 
    must be constant.  

    Recall that $\overline A$ is the smallest 
    closed set containing $A$, equivalently it is 
    the intersection of all closed 
    subsets of $X$ containing $A$.  
    Denote $f(A) = \{c_A\}$. 

    Let $f : \overline A \to \{0,1\}$ be some continuous 
    function. 

    Suppose, $x,y \in \overline A$ with 
    \[
        f(x) = c_x, \ f(y) = c_y. 
    \] 
    
    The fiber of $c_x$ is an open set $V_x$ of $\overline A$ 
    containing $x$.     
    Likewise, the fiber of $c_y$ is an open set
    $V_y$ of $\overline A$ containing $y$. 
    
    Then $f(V_x \cap A) = \{c_x\}$ and 
    $f(V_y \cap A) = \{c_y\}$, but since 
    these are subsets of $A$, and 
    \[
        \im f\bigm|_A = \{c_A\},
    \]
    we must have 
    \[
        c_x = c_A = c_y. 
    \]

    Therefore the space is connected. 

    Note: I should have been more rigorous 
    in explaining that we can do this by arguing that 
    since the points were in the closure, any open 
    neighborhood necessarily intersects $A$. 
    
\end{proof}

\begin{prob}
    Let $C$ be a connected component of a topological 
    space $X$. Prove that $C$ is closed in $X$. 

    You may use $A$ connected 
    $\Rightarrow$ $\overline A$ connected. 
\end{prob}
\begin{proof}
    $C$ is a maximal connected subset. 

    Thus there does not exist any 
    $D \subseteq X$ containing $C$ which is also connected. 
    
    In particular, since $C \subseteq \overline C$ 
    we must have that 
    \[
        C = \overline C. 
    \]
    Otherwise, $\overline C$ would be a larger connected 
    subspace ($C$ connected $\Rightarrow$ 
    $\overline C$ connected) containing all of $C$,  
    contradicting the maximality of $C$. 
\end{proof}

\begin{prob}
    Let $X$ be locally connected. Let $U \subseteq X$ 
    be open, and let $C$ be a connected component 
    of $U$. 

    Prove that $C$ is open in $X$. 
\end{prob}
\begin{proof}
    Since $X$ is locally connected, then for every  
    $x \in X$ we can, for every open 
    neighborhood $U_x \supseteq x$, find an open 
    connected set $V_x \supseteq x$ contained in $U_x$. 

    We suppose $U$ is an open set in $X$ and 
    that $C$ is a maximal connected subset of $U$. 

    For any $x \in C$, we can find a connected open 
    neighborhood 
    $V_x$ with $V_x \subseteq U$. 

    Since $C$ is maximal, $x \in C$, and $V_x$ is connected 
    we have $V_x \subseteq C$ for every $x$ (might 
    need one more step to be completely 
    justified). 
    Then 
    \[
        \bigcup_{x} V_x = C. 
    \]
    Since open sets are closed under arbitrary unions, 
    $C$ is open. 

\end{proof}



\end{document}
